\documentclass[11pt]{article}

\usepackage[a4paper,margin=1in]{geometry}
\usepackage[T1]{fontenc}
\usepackage[utf8]{inputenc}
\usepackage{lmodern}
\usepackage{microtype}
\usepackage{amsmath,amssymb,amsthm,mathtools}
\usepackage{enumitem}
\usepackage[colorlinks=true,linkcolor=blue,citecolor=blue,urlcolor=blue]{hyperref}

\title{Substitutability, Nonlinear Jacobi, and Sinkhorn as a Special Case}
\author{A short pedagogical note}
\date{\today}

\newcommand{\R}{\mathbb{R}}
\newcommand{\one}{\mathbf{1}}
\newcommand{\dd}{\mathrm{d}}
\newcommand{\diag}{\operatorname{diag}}
\newcommand{\supp}{\operatorname{supp}}
\newcommand{\e}{\mathrm{e}}
\newcommand{\abs}[1]{\left\lvert #1\right\rvert}
\newcommand{\braces}[1]{\left\{#1\right\}}
\newcommand{\ip}[2]{\left\langle #1,#2\right\rangle}

\theoremstyle{definition}
\newtheorem{definition}{Definition}[section]
\newtheorem{assumption}[definition]{Assumption}
\newtheorem{example}[definition]{Example}
\newtheorem{remark}[definition]{Remark}

\theoremstyle{plain}
\newtheorem{proposition}[definition]{Proposition}
\newtheorem{theorem}[definition]{Theorem}
\newtheorem{lemma}[definition]{Lemma}

\begin{document}
\maketitle

\begin{abstract}
Sinkhorn's algorithm is usually introduced as a variational algorithm: alternating Bregman projections, or coordinate ascent/descent on the dual of entropically regularized optimal transport. There is another point of view, emphasized in recent work around Galichon's equilibrium-flow program: Sinkhorn is a special case of a nonlinear Jacobi market-clearing algorithm. This second point of view does not require the equilibrium equations to be the gradient of a convex potential. The relevant structure is substitutability, formalized through \emph{Z-functions}, \emph{M-functions}, and \emph{$M_0$-functions}. A familiar linear sufficient condition is a nonsingular M-matrix condition; in sign-uncertain formulations this is often expressed through an H-matrix condition. A worked lossy-transport example illustrates how the same Sinkhorn-like scaling formulas can remain convergent even when no convex OT potential exists.
\end{abstract}

\tableofcontents

\section{The basic fixed-point problem}

Let $Z$ be a finite set of coordinates, and let
\[
        Q : D \subseteq \R^Z \longrightarrow \R^Z
\]
be a continuous map. In economic applications, $p\in \R^Z$ is a price vector and $Q_z(p)$ is excess supply in market $z$. The equilibrium problem is
\begin{equation}\label{eq:Qzero}
        Q(p)=0.
\end{equation}

The nonlinear Jacobi update clears each coordinate separately, holding the other coordinates fixed at their old values. Formally, for $p^t\in D$, define $p^{t+1}=T(p^t)$ by
\begin{equation}\label{eq:jacobi}
        Q_z\bigl(p^{t+1}_z,p^t_{-z}\bigr)=0,\qquad z\in Z,
\end{equation}
where $p_{-z}$ denotes all coordinates except $z$.

\begin{assumption}[Responsiveness]
For every $z$ and every $p_{-z}$, the one-dimensional equation
\[
        Q_z(\pi,p_{-z})=0
\]
has at least one solution in $\pi$. When there is more than one solution, we select the smallest one:
\[
        T_z(p)=\min\braces{\pi:Q_z(\pi,p_{-z})=0}.
\]
\end{assumption}

\begin{remark}[Jacobi versus Gauss--Seidel]
Equation \eqref{eq:jacobi} is the parallel Jacobi version. The familiar implementation of Sinkhorn often uses a Gauss--Seidel order: update rows using the old columns, then update columns using the new rows. Both are coordinate-clearing schemes; the order-theoretic hypotheses below are most transparent for Jacobi, and analogous monotone arguments apply to ordered coordinate updates.
\end{remark}

\section{Subsolutions and the first monotonicity argument}

\begin{definition}[Subsolution and supersolution]
A vector $p$ is a subsolution if $Q(p)\leq 0$ componentwise. It is a supersolution if $Q(p)\geq 0$.
\end{definition}

The intended interpretation is that subsolution prices are too low and supersolution prices are too high. For this interpretation to be valid, own-price effects and cross-price effects need the right signs.

\begin{definition}[Diagonal isotonicity and Z-functions]
The map $Q$ is diagonally isotone if, for each $z$, $Q_z(p_z,p_{-z})$ is nondecreasing in $p_z$.

It is a Z-function if, for each $z$, $Q_z(p_z,p_{-z})$ is antitone in $p_{-z}$: whenever $p_{-z}\leq p'_{-z}$,
\[
        Q_z(p_z,p_{-z})\geq Q_z(p_z,p'_{-z}).
\]
For $C^1$ maps, this means
\[
        \frac{\partial Q_z}{\partial p_{z'}}(p)\leq 0,\qquad z'\neq z.
\]
\end{definition}

\begin{proposition}[Monotone preservation of subsolutions]\label{prop:sub}
Assume $Q$ is diagonally isotone and a Z-function. If $p^t$ is a subsolution, then
\[
        p^t\leq p^{t+1}\quad\text{and}\quad p^{t+1}\text{ is again a subsolution}.
\]
The analogous statement holds for supersolutions, with inequalities reversed.
\end{proposition}

\begin{proof}
Since $p^t$ is a subsolution,
\[
        Q_z(p_z^t,p^t_{-z})\leq 0=Q_z(p_z^{t+1},p^t_{-z}).
\]
By diagonal isotonicity, $p_z^t\leq p_z^{t+1}$ for every $z$. Hence $p^t\leq p^{t+1}$. Now, using the Z-property,
\[
        Q_z(p_z^{t+1},p_{-z}^{t+1})\leq Q_z(p_z^{t+1},p_{-z}^{t})=0.
\]
Thus $p^{t+1}$ is a subsolution.
\end{proof}

The preceding proposition is useful but not sufficient for convergence.

\begin{example}[Why Z-structure alone is not enough]
Let
\[
        Q(p)=\begin{pmatrix}1&-2\\-2&1\end{pmatrix}p.
\]
Then $Q$ is diagonally isotone and a Z-function, but the Jacobi update is
\[
        p_1^{t+1}=2p_2^t,
        \qquad
        p_2^{t+1}=2p_1^t.
\]
Except from the fixed point $p=0$, the iterates diverge. The problem is that cross-effects dominate own-effects.
\end{example}

\section{Nonreversingness, M-functions, and inverse isotonicity}

The missing condition says that an increase in all prices cannot make all excess supplies fall. This is the nonlinear counterpart of the linear M-matrix idea.

\begin{definition}[Nonreversingness]
Let $Q:D\to\R^Z$.
\begin{enumerate}[label=(\roman*)]
\item $Q$ is strongly nonreversing if
\[
        p\geq p',\quad Q(p)\leq Q(p')
        \quad\Longrightarrow\quad p=p'.
\]
\item $Q$ is weakly nonreversing if
\[
        p\geq p',\quad Q(p)\leq Q(p')
        \quad\Longrightarrow\quad Q(p)=Q(p').
\]
\end{enumerate}
\end{definition}

\begin{definition}[M- and $M_0$-functions]
An M-function is a Z-function that is strongly nonreversing. An $M_0$-function is a Z-function that is weakly nonreversing.
\end{definition}

Thus an M-function has two ingredients:
\begin{enumerate}[label=(\alph*)]
\item substitutes: increasing other prices lowers excess supply for a given good;
\item own-price dominance: substitution effects cannot reverse the direction of the total price movement.
\end{enumerate}

\begin{theorem}[Inverse isotonicity]\label{thm:inverse}
Let $Q$ be a Z-function. Then:
\begin{enumerate}[label=(\roman*)]
\item $Q$ is an M-function if and only if it is inverse isotone:
\[
        Q(p)\leq Q(p')\quad\Longrightarrow\quad p\leq p'.
\]
\item $Q$ is an $M_0$-function if and only if its inverse is isotone as a set-valued map. In the single-valued notation, a convenient equivalent statement is
\[
        Q(p)\leq Q(p')
        \quad\Longrightarrow\quad
        Q(p\wedge p')=Q(p),\qquad Q(p\vee p')=Q(p'),
\]
where $p\wedge p'$ and $p\vee p'$ denote componentwise minimum and maximum.
\end{enumerate}
\end{theorem}

This theorem is the order-theoretic reason the algorithm works. It replaces convexity of a potential by monotonicity of the inverse equilibrium map.

\section{Convergence of nonlinear Jacobi}

\begin{theorem}[Global monotone convergence for M-functions]\label{thm:Mconv}
Assume $Q$ is continuous, responsive, diagonally isotone, and an M-function. Assume also that there exist a subsolution $\underline p$ and a supersolution $\overline p$. Then the equation $Q(p)=0$ has a unique solution $p^*$, and Jacobi iterates starting from any subsolution or any supersolution converge monotonically to $p^*$.
\end{theorem}

\begin{proof}[Proof sketch]
Starting from a subsolution, Proposition \ref{prop:sub} gives an increasing sequence of subsolutions. Starting from a supersolution gives a decreasing sequence of supersolutions. By inverse isotonicity, every subsolution lies below every supersolution: indeed, if $Q(\underline p)\leq0\leq Q(\overline p)$, then Theorem \ref{thm:inverse} yields $\underline p\leq\overline p$. Hence the lower and upper Jacobi sequences are monotone and bounded. Their limits solve $Q(p)=0$ by continuity and responsiveness. If two solutions $p,p'$ existed, then both $Q(p)\leq Q(p')$ and $Q(p')\leq Q(p)$; inverse isotonicity gives $p\leq p'$ and $p'\leq p$, hence $p=p'$.
\end{proof}

\begin{theorem}[The $M_0$ case]
If $Q$ is an $M_0$-function, uniqueness may fail. Under the same continuity, responsiveness, and lower-upper solution hypotheses, the solution set is a lattice. Jacobi started from the meet or the join of a subsolution and a supersolution converges to an extremal solution.
\end{theorem}

\begin{remark}[Why $M_0$ matters for Sinkhorn]
Balanced optimal transport is invariant under adding the same constant to all dual potentials of one sign convention. Therefore the price-clearing map is typically $M_0$, not M, before a normalization is imposed. Fixing one potential, or imposing a gauge such as $\sum p_z=0$, removes the null direction and turns the reduced system into an M-system under the usual connectivity assumptions.
\end{remark}

\section{A differentiable M-matrix/H-matrix certificate}

Suppose $Q$ is $C^1$ on an order-convex domain, and write $J(p)=DQ(p)$.

\begin{definition}[Z-, M-, and H-matrices]
A matrix $A$ is a Z-matrix if $A_{ij}\leq0$ for $i\neq j$. A nonsingular M-matrix is a Z-matrix whose inverse is nonnegative; equivalently, $A=sI-B$ with $B\geq0$ and $s>\rho(B)$.

The comparison matrix of $A$ is
\[
        \langle A\rangle_{ii}=|A_{ii}|,
        \qquad
        \langle A\rangle_{ij}=-|A_{ij}|\quad(i\neq j).
\]
The matrix $A$ is an H-matrix if $\langle A\rangle$ is an M-matrix. If $A$ has positive diagonal entries, one often says $A$ is an $H_+$-matrix. For Z-matrices with positive diagonal, the $H_+$ and nonsingular M-matrix conditions coincide.
\end{definition}

The following criterion is often the easiest way to recognize the condition informally remembered as an ``H-matrix'' condition.

\begin{proposition}[Weighted diagonal dominance implies an M-function]\label{prop:H}
Assume $Q$ is $C^1$ and satisfies
\begin{equation}\label{eq:zpartials}
        \frac{\partial Q_i}{\partial p_j}(p)\leq0,
        \qquad i\neq j,
\end{equation}
for all $p$. Suppose there exists a vector $\lambda\gg0$ such that
\begin{equation}\label{eq:weightedDD}
        \lambda^\top J(p)\gg0
        \qquad\text{for every }p.
\end{equation}
Equivalently, for each column $j$,
\begin{equation}\label{eq:columnDD}
        \lambda_j\,\frac{\partial Q_j}{\partial p_j}(p)
        >
        \sum_{i\neq j}\lambda_i\left|\frac{\partial Q_i}{\partial p_j}(p)\right|.
\end{equation}
Then $Q$ is an M-function. If the strict inequalities are weakened to non-strict inequalities and $\lambda^\top J(p)\geq0$, then $Q$ is an $M_0$-function.
\end{proposition}

\begin{proof}
Condition \eqref{eq:zpartials} says that $Q$ is a Z-function. Let $p\geq p'$ and set $\delta=p-p'\geq0$. If $\delta\neq0$, then
\[
        \lambda^\top\bigl(Q(p)-Q(p')\bigr)
        =\int_0^1 \lambda^\top J(p'+t\delta)\,\delta\,\dd t
        >0
\]
by \eqref{eq:weightedDD}. Hence it is impossible to have $Q(p)\leq Q(p')$, because then the left side would be nonpositive. This proves strong nonreversingness. The weak version is the same argument with $\geq0$: if $Q(p)-Q(p')\leq0$ and its positive weighted sum is nonnegative, then the vector must be zero.
\end{proof}

\begin{remark}[Relation to H-matrices]
When $J(p)$ is a Z-matrix, strict weighted diagonal dominance such as \eqref{eq:columnDD} implies that $J(p)$ is a nonsingular M-matrix, hence an $H_+$-matrix. The H-matrix language is useful in numerical analysis because it survives some sign changes. The economic/substitution theory usually keeps the Z-sign pattern explicit and then speaks of M-functions.
\end{remark}

\section{Sinkhorn as the canonical special case}

Let $X$ and $Y$ be finite sets. Let $a=(a_x)_{x\in X}$ and $b=(b_y)_{y\in Y}$ be positive masses with equal total mass, and let $K_{xy}>0$ be a positive kernel. Entropically regularized optimal transport seeks a coupling $\pi$ with row sums $a$ and column sums $b$ of the form
\begin{equation}\label{eq:scalingform}
        \pi_{xy}=r_x K_{xy}s_y,
        \qquad r_x>0,\qquad s_y>0.
\end{equation}
The row and column constraints are
\[
        r_x\sum_y K_{xy}s_y=a_x,
        \qquad
        s_y\sum_x K_{xy}r_x=b_y.
\]
Sinkhorn's scaling updates are
\begin{equation}\label{eq:sinkhorn-rs}
        r_x^{t+1}=\frac{a_x}{\sum_y K_{xy}s_y^t},
        \qquad
        s_y^{t+1}=\frac{b_y}{\sum_x K_{xy}r_x^t}.
\end{equation}
The common alternating implementation uses the new $r^{t+1}$ in the second equation; \eqref{eq:sinkhorn-rs} is the parallel Jacobi version.

To put this into the M-function notation, set
\[
        u_x=\log r_x,
        \qquad
        v_y=-\log s_y,
        \qquad
        \pi_{xy}(u,v)=K_{xy}\e^{u_x-v_y}.
\]
Define $Q:\R^X\times\R^Y\to\R^X\times\R^Y$ by
\begin{align}
        Q_x(u,v) &= \sum_{y\in Y}\pi_{xy}(u,v)-a_x,\label{eq:Qsinkx}\\
        Q_y(u,v) &= b_y-\sum_{x\in X}\pi_{xy}(u,v).\label{eq:Qsinky}
\end{align}
Then $Q(u,v)=0$ is exactly the row-column marginal system. The Jacobi equations $Q_x(u_x^{t+1},v^t)=0$ and $Q_y(u^t,v_y^{t+1})=0$ give
\begin{align*}
        u_x^{t+1}
        &=\log a_x-\log\sum_yK_{xy}\e^{-v_y^t},\\
        v_y^{t+1}
        &=\log\sum_xK_{xy}\e^{u_x^t}-\log b_y.
\end{align*}
Exponentiating gives exactly \eqref{eq:sinkhorn-rs}.

\subsection*{The Z- and $M_0$-structure}
The Jacobian of \eqref{eq:Qsinkx}--\eqref{eq:Qsinky} has entries
\begin{align*}
        \frac{\partial Q_x}{\partial u_x}&=\sum_y\pi_{xy},
        &\frac{\partial Q_x}{\partial v_y}&=-\pi_{xy},\\
        \frac{\partial Q_y}{\partial v_y}&=\sum_x\pi_{xy},
        &\frac{\partial Q_y}{\partial u_x}&=-\pi_{xy}.
\end{align*}
All off-diagonal derivatives are nonpositive, so $Q$ is a Z-function. Moreover, every column sum of $J=DQ$ is zero:
\[
        \one^\top J=0.
\]
Therefore the map is weakly nonreversing but not strongly nonreversing. The missing strong condition is exactly the additive gauge invariance
\[
        (u,v)\longmapsto (u+c\one_X,v+c\one_Y),
\]
under which $u_x-v_y$ and hence $\pi_{xy}$ are unchanged. Thus balanced Sinkhorn is naturally an $M_0$ example. After fixing one potential or imposing a normalization, the reduced Jacobian is a weighted graph Laplacian minor; if the support graph of $K$ is connected, that minor is a nonsingular M-matrix.

\subsection*{Why Sinkhorn is also variational}
For entropic OT there is additionally a convex dual potential. In one sign convention,
\[
        \Phi(u,v)
        =\sum_{x,y}K_{xy}\e^{u_x-v_y}
          -\sum_xa_xu_x+\sum_yb_yv_y,
\]
up to irrelevant constants and sign choices. The first-order conditions are precisely $Q(u,v)=0$. Sinkhorn can therefore be interpreted as coordinate descent/ascent on a convex objective. The order-theoretic point is different: the coordinate update only needs the market-clearing equations $Q_z=0$ and the M- or $M_0$-structure. It does not need $Q$ to be a gradient.


\section{A concrete OT-inspired non-variational example}\label{sec:lossy}

The preceding Sinkhorn example is perfectly conservative: one unit leaving an origin arrives as one unit at a destination. Here is a small modification which keeps the same transport geometry and the same scaling form, but is generally not variational.

Let $X$ and $Y$ be finite sets, let $K_{xy}>0$, and introduce positive outside-option coefficients $\sigma_x,\tau_y>0$. We write
\[
        r_x=\e^{u_x},\qquad s_y=\e^{-v_y},\qquad
        \pi_{xy}=r_xK_{xy}s_y .
\]
Think of $\pi_{xy}$ as the physical amount shipped from $x$ to $y$. Suppose, however, that only a fraction $\eta_{xy}\in(0,1]$ is effective at the destination. The equilibrium equations are
\begin{align}
        \sigma_x r_x+\sum_{y\in Y}\pi_{xy}&=a_x,
        &&x\in X,\label{eq:lossy-row}\\
        \tau_y s_y+\sum_{x\in X}\eta_{xy}\pi_{xy}&=b_y,
        &&y\in Y.\label{eq:lossy-col}
\end{align}
The first equation says that the mass available at origin $x$ is split between the transport network and an outside sink. The second says that destination $y$ can be served either by effective arrivals or by a local outside option. This is not meant to be a standard OT model; it is a deliberately simple OT-shaped market-clearing system.

In the signed logarithmic variables $(u,v)$ define
\begin{align*}
        Q_x(u,v)
        &=\sigma_x\e^{u_x}+\sum_y K_{xy}\e^{u_x-v_y}-a_x,\\
        Q_y(u,v)
        &=b_y-\tau_y\e^{-v_y}-\sum_x\eta_{xy}K_{xy}\e^{u_x-v_y}.
\end{align*}
Then $Q=0$ is exactly \eqref{eq:lossy-row}--\eqref{eq:lossy-col}. Clearing one coordinate at a time gives the Sinkhorn-like update
\begin{equation}\label{eq:lossy-update}
        r_x^{t+1}
        =\frac{a_x}{\sigma_x+\sum_yK_{xy}s_y^t},
        \qquad
        s_y^{t+1}
        =\frac{b_y}{\tau_y+\sum_x\eta_{xy}K_{xy}r_x^t}.
\end{equation}
Thus the algorithm still alternately rescales rows and columns. The only difference from Sinkhorn is that the marginal equations contain outside terms and the column equation sees the distorted arrivals $\eta_{xy}\pi_{xy}$.

\subsection*{The M-matrix certificate}

The Jacobian has the same Z-sign pattern as Sinkhorn:
\begin{align*}
        \frac{\partial Q_x}{\partial u_x}
        &=\sigma_xr_x+\sum_y\pi_{xy},
        &\frac{\partial Q_x}{\partial v_y}
        &=-\pi_{xy},\\
        \frac{\partial Q_y}{\partial u_x}
        &=-\eta_{xy}\pi_{xy},
        &\frac{\partial Q_y}{\partial v_y}
        &=\tau_ys_y+\sum_x\eta_{xy}\pi_{xy}.
\end{align*}
Assume $0<\eta_{xy}\leq1$ and
\begin{equation}\label{eq:lossy-condition}
        \tau_y>
        \sum_x(1-\eta_{xy})K_{xy}\frac{a_x}{\sigma_x},
        \qquad y\in Y.
\end{equation}
After one update, $r_x\leq a_x/\sigma_x$. On the invariant domain where these inequalities hold, the vector $\lambda=\one$ certifies weighted diagonal dominance. Indeed, the column sum corresponding to $u_x$ is
\[
        \sigma_xr_x+
        \sum_y(1-\eta_{xy})\pi_{xy}>0,
\]
while the column sum corresponding to $v_y$ is
\[
        \tau_ys_y-
        \sum_x(1-\eta_{xy})\pi_{xy}
        =s_y\left[\tau_y-
        \sum_x(1-\eta_{xy})K_{xy}r_x\right]>0
\]
by \eqref{eq:lossy-condition}. Hence $\one^\top DQ\gg0$ on the invariant domain. By Proposition~\ref{prop:H}, this is an M-function there, so the coordinate-clearing iteration converges monotonically once initialized from a subsolution or supersolution in that domain.

The system is generally not variational. If $Q=\nabla\Phi$ for a $C^2$ potential $\Phi$, then the cross-partials would have to agree:
\[
        \frac{\partial Q_x}{\partial v_y}
        =\frac{\partial Q_y}{\partial u_x}.
\]
Here this would require $-\pi_{xy}=-\eta_{xy}\pi_{xy}$ for every active arc, hence $\eta_{xy}=1$. Losses with $\eta_{xy}\neq1$ therefore break exactness of the one-form $\sum_z Q_z\,\dd p_z$, even though the M-structure and the Sinkhorn-shaped update remain.

\subsection*{Numerical instance}

Take two origins and two destinations,
\[
        K=\begin{pmatrix}
        1&\e^{-1}\\
        \e^{-1}&1
        \end{pmatrix},\qquad
        a=b=\begin{pmatrix}1\\1\end{pmatrix},
        \qquad
        \sigma=\tau=\begin{pmatrix}1\\1\end{pmatrix},
\]
with loss factors
\[
        \eta=\begin{pmatrix}
        0.80&0.90\\
        0.70&0.85
        \end{pmatrix}.
\]
Condition \eqref{eq:lossy-condition} is immediate:
\[
        0.20+0.30\e^{-1}\simeq0.310<1,
        \qquad
        0.10\e^{-1}+0.15\simeq0.187<1.
\]
Starting from $r^0=s^0=(1,1)$, the parallel update \eqref{eq:lossy-update} gives
\[
\begin{array}{c|cccc}
        t&r_1^t&r_2^t&s_1^t&s_2^t\\
        \hline
        0&1.0000&1.0000&1.0000&1.0000\\
        1&0.4223&0.4223&0.4860&0.4585\\
        2&0.6043&0.6108&0.6913&0.6672\\
        3&0.5163&0.5204&0.6095&0.5817\\
        6&0.5391&0.5441&0.6388&0.6119
\end{array}
\]
The limit is approximately
\[
        r^*=(0.5371,0.5421),
        \qquad
        s^*=(0.6372,0.6103).
\]
The corresponding transport matrix is
\[
        \pi^*\simeq
        \begin{pmatrix}
        0.3423&0.1206\\
        0.1271&0.3308
        \end{pmatrix}.
\]
One can check the clearing equations directly:
\begin{align*}
        r_1^*+\pi_{11}^*+\pi_{12}^*&\simeq1,
        &r_2^*+\pi_{21}^*+\pi_{22}^*&\simeq1,\\
        s_1^*+0.80\pi_{11}^*+0.70\pi_{21}^*&\simeq1,
        &s_2^*+0.90\pi_{12}^*+0.85\pi_{22}^*&\simeq1.
\end{align*}
This is the closest finite-dimensional analogue of Sinkhorn in the note: positive kernel, multiplicative scalings, marginal-clearing equations, and closed-form row/column updates. The reason it converges, however, is the M-function condition, not a convex OT potential.

\section{A minimal algebraic non-variational example}

Let $a,b>0$ and $0\leq\lambda,\mu<1$. Consider
\begin{align*}
        Q_1(p_1,p_2)&=\e^{p_1}-\lambda\e^{p_2}-a,\\
        Q_2(p_1,p_2)&=\e^{p_2}-\mu\e^{p_1}-b.
\end{align*}
Then
\[
        J(p)=
        \begin{pmatrix}
        \e^{p_1} & -\lambda\e^{p_2}\\
        -\mu\e^{p_1} & \e^{p_2}
        \end{pmatrix},
        \qquad
        \one^\top J(p)=\bigl((1-\mu)\e^{p_1},(1-\lambda)\e^{p_2}\bigr)\gg0.
\]
Thus Proposition \ref{prop:H} applies: $Q$ is an M-function. However, unless the cross-partials happen to agree everywhere,
\[
        \frac{\partial Q_1}{\partial p_2}=-\lambda\e^{p_2}
        \neq
        -\mu\e^{p_1}=\frac{\partial Q_2}{\partial p_1},
\]
so $Q$ is not the gradient of any scalar potential. The Jacobi update is nevertheless simple and convergent:
\[
        p_1^{t+1}=\log\bigl(a+\lambda\e^{p_2^t}\bigr),
        \qquad
        p_2^{t+1}=\log\bigl(b+\mu\e^{p_1^t}\bigr).
\]
The unique solution is obtained by solving the linear equations in $x=\e^{p_1}$ and $y=\e^{p_2}$:
\[
        x=a+\lambda y,
        \qquad y=b+\mu x,
\]
so
\[
        x^*=\frac{a+\lambda b}{1-\lambda\mu},
        \qquad
        y^*=\frac{b+\mu a}{1-\lambda\mu}.
\]
This example is deliberately elementary, but it isolates the important point: the Jacobi/Sinkhorn-type update can converge for structural reasons that are not variational.

\section{Equilibrium flows and imperfect transferability}

The broader Galichon--Samuelson--Vernet framework replaces the additive OT inequality
\[
        p_x-p_y\geq -c_{xy}
\]
by a nonlinear arc condition
\[
        p_x\geq G_{xy}(p_y),
\]
where $G_{xy}$ is increasing. An equilibrium flow consists of a node price vector $p$, nonnegative arc flows $\mu_{xy}$, and net node flows $q$ such that
\begin{align*}
        \nabla^\top\mu&=q,\\
        p_x&\geq G_{xy}(p_y),\qquad xy\in A,\\
        \mu_{xy}>0&\Longrightarrow p_x=G_{xy}(p_y).
\end{align*}
The additive case $G_{xy}(p_y)=p_y-c_{xy}$ is the usual min-cost-flow/transport duality. Nonlinear $G_{xy}$ accommodates time-dependent routing, nonlinear taxes, and imperfectly transferable utility in matching. In these cases, there need not be a convex objective whose gradient is the equilibrium equation. What remains is substitutability: increasing the relevant price at one node makes neighboring opportunities less attractive in the appropriate signed coordinates. The resulting equilibrium correspondence has inverse-isotone/lattice structure, and coordinate-clearing price updates are the natural generalization of Sinkhorn.

\section{Takeaway}

Sinkhorn sits at the intersection of two structures:
\begin{enumerate}[label=(\arabic*)]
\item a variational structure: entropic OT is a convex optimization problem;
\item an order structure: the marginal equations form an $M_0$-system, and Sinkhorn is a coordinate-clearing Jacobi/Gauss--Seidel method.
\end{enumerate}
Galichon's more general viewpoint keeps the second structure and drops the first. The lossy-transport example in Section~\ref{sec:lossy} makes this separation concrete: the update has the same multiplicative scaling form as Sinkhorn, but asymmetric effective arrivals destroy the cross-partial symmetry required by a potential. The matrix condition to remember is best described as a nonlinear M-function condition. In smooth finite-dimensional problems, a practical certificate is a Z-Jacobian plus weighted diagonal dominance, equivalently an M-matrix or $H_+$-matrix condition after the relevant sign convention and normalization.

\begin{thebibliography}{9}

\bibitem{GalichonJacquet2024}
A. Galichon and A. Jacquet.
\newblock \emph{Substitutability, equilibrium transport, and matching models}.
\newblock arXiv:2405.07628, 2024.

\bibitem{GSV2025}
A. Galichon, L. Samuelson, and L. Vernet.
\newblock Unified gross substitutes and inverse isotonicity for equilibrium problems.
\newblock \emph{Theoretical Economics}, 20(3):941--971, 2025.

\bibitem{GSV2022flows}
A. Galichon, L. Samuelson, and L. Vernet.
\newblock \emph{The existence of equilibrium flows}.
\newblock arXiv:2209.04426, 2022.

\bibitem{OrtegaRheinboldt1970}
J. M. Ortega and W. C. Rheinboldt.
\newblock \emph{Iterative Solution of Nonlinear Equations in Several Variables}.
\newblock Academic Press, 1970.

\bibitem{SinkhornKnopp1967}
R. Sinkhorn and P. Knopp.
\newblock Concerning nonnegative matrices and doubly stochastic matrices.
\newblock \emph{Pacific Journal of Mathematics}, 21(2):343--348, 1967.

\bibitem{QuGalichonGaoUgander2025}
Z. Qu, A. Galichon, W. Gao, and J. Ugander.
\newblock On Sinkhorn's algorithm and choice modeling.
\newblock \emph{Operations Research}, 2025.

\end{thebibliography}

\end{document}
